\documentclass[../cours_1re.tex]{subfiles}
\begin{document}
\setcounter{chapter}{14}
\chapter{Identités et fonctions trigonométriques}
\label{ch:15}
\begin{objectif}
Approfondissement hors du noyau du programme : passer du cercle aux identités, puis aux fonctions périodiques. Prérequis : chapitre 14 ; le produit scalaire du chapitre 16 pour la première preuve.
\end{objectif}
\section{Des angles aux identités}
Les identités suivantes sont utiles, mais ne remplacent pas le dessin du cercle.
Pour tous réels $a,b$,
\[
\cos(a-b)=\cos a\cos b+\sin a\sin b,\qquad
\sin(a+b)=\sin a\cos b+\cos a\sin b.
\]
\begin{demonstration}
Les vecteurs unitaires faisant les angles $a$ et $b$ ont pour coordonnées
$(\cos a,\sin a)$ et $(\cos b,\sin b)$. Leur produit scalaire vaut à la fois
$\cos(a-b)$ et $\cos a\cos b+\sin a\sin b$ (chapitre 16, dont la preuve ne suppose
aucune formule d'addition). Remplacer $b$ par $-b$ donne
$\cos(a+b)=\cos a\cos b-\sin a\sin b$.
Enfin $\sin t=\cos(\pi/2-t)$ et les angles associés du chapitre 14 donnent
la formule du sinus. Les preuves sont donc indépendantes de la dérivation du sinus.
\end{demonstration}
On obtient également $\sin(a-b)=\sin a\cos b-\cos a\sin b$.
Avec $a=b=x$ :
\[
\sin(2x)=2\sin x\cos x,\qquad
\cos(2x)=\cos^2x-\sin^2x=1-2\sin^2x=2\cos^2x-1.
\]
Ces égalités donnent $\sin^2x=(1-\cos2x)/2$ et
$\cos^2x=(1+\cos2x)/2$. Pour transformer un produit en somme, ajouter
les formules de $\cos(a-b)$ et $\cos(a+b)$ :
\[
2\cos a\cos b=\cos(a-b)+\cos(a+b).
\]
\begin{exemple}
$\cos(\pi/12)=\cos(\pi/4-\pi/6)=(\sqrt6+\sqrt2)/4$.
La valeur est positive car l'angle est dans le premier quadrant.
\end{exemple}

\section{Pourquoi le radian est indispensable à la dérivation}
Pour $0<h<\pi/2$, la comparaison des aires du triangle inscrit,
du secteur d'angle $h$ et du triangle tangent au cercle unité donne
\[
\tfrac12\sin h<\tfrac12h<\tfrac12\tan h,\qquad
\cos h<\frac{\sin h}{h}<1.
\]
Ici l'aire du secteur est $h/2$, par proportionnalité à l'angle :
un disque d'aire $\pi$ correspond à $2\pi$ radians. On utilise ce résultat
géométrique d'aire, qui peut être construit par approximation polygonale.
Comme $0\le1-\cos h=2\sin^2(h/2)\le h^2/2$,
$\cos h\to1$. L'encadrement, puis la parité, donnent
\[
\lim_{h\to0}\frac{\sin h}{h}=1,\qquad
\lim_{h\to0}\frac{\cos h-1}{h}=0.
\]
Les bornes sur $\sin h$ utilisées ici se prolongent aux petits $h$ négatifs
par imparité. Elles prouvent aussi la continuité des fonctions trigonométriques,
en utilisant les formules d'addition.
\begin{demonstration}[dérivées]
La formule d'addition donne
\[
\frac{\sin(x+h)-\sin x}{h}
=\sin x\,\frac{\cos h-1}{h}+\cos x\,\frac{\sin h}{h}\longrightarrow\cos x.
\]
Le même calcul donne $(\cos x)'=-\sin x$. Pour $\cos x\ne0$,
la règle du quotient donne $(\tan x)'=1/\cos^2x$.
\end{demonstration}
\begin{attention}
Ces formules supposent des angles en radians. Pour $g(t)=\sin(\pi t/180)$,
où $t$ est une mesure en degrés, $g'(t)=(\pi/180)\cos(\pi t/180)$.
\end{attention}

\section{Lire et construire des oscillations}
Sur $[0,2\pi]$, le sinus croît sur $[0,\pi/2]$, décroît sur
$[\pi/2,3\pi/2]$, puis croît ; c'est le signe du cosinus qui le prouve.
La périodicité prolonge cette description à $\R$.
\begin{center}
\begin{tikzpicture}
\begin{axis}[width=.92\linewidth,height=5cm,axis lines=middle,
xmin=0,xmax=6.6,ymin=-1.25,ymax=1.25,
xtick={0,1.5708,3.1416,4.7124,6.2832},
xticklabels={$0$,$\pi/2$,$\pi$,$3\pi/2$,$2\pi$},ytick={-1,1},
legend style={at={(.98,.98)},anchor=north east}]
\addplot[accent,thick,domain=0:6.2832,samples=150]{sin(deg(x))};
\addlegendentry{$\sin x$}
\addplot[black,dashed,domain=0:6.2832,samples=150]{cos(deg(x))};
\addlegendentry{$\cos x$}
\end{axis}
\end{tikzpicture}
\end{center}
Une oscillation $A\cos(\omega t+\varphi)$, avec $A>0$ et $\omega>0$,
a pour amplitude $A$ et période $T=2\pi/\omega$ : augmenter $t$ de $T$
ajoute $2\pi$ à l'angle. Sa dérivée seconde vaut $-\omega^2$ fois la fonction.
Elle résout donc $y''+\omega^2y=0$, modèle idéal d'un ressort sans frottement.

\section{Rotation et conservation des longueurs}
L'application
\[
(x,y)\longmapsto(x\cos\theta-y\sin\theta,\ x\sin\theta+y\cos\theta)
\]
fait tourner les vecteurs d'angle $\theta$. Pour un vecteur de longueur $r$
et d'angle $a$, les formules d'addition transforment ses coordonnées en
$(r\cos(a+\theta),r\sin(a+\theta))$ ; le vecteur nul est fixe.
En développant la somme des carrés des coordonnées, les termes croisés
s'annulent : la longueur est conservée. Composer deux rotations revient
à additionner les angles. La rotation inverse est celle d'angle $-\theta$.

\begin{finchapitre}
\section{Utilité et prolongements}
Les sinus et cosinus décrivent des positions et des vitesses périodiques.
Les produits transformés en sommes serviront au calcul d'intégrales.

En licence, les nombres complexes écrivent la rotation comme une multiplication.
Les séries de Fourier décomposent ensuite un signal en oscillations ; leur convergence
demande des outils d'analyse supplémentaires, et ne découle pas d'un simple dessin.
\begin{retenir}
Les identités viennent de la géométrie du cercle. Le radian rend
$(\sin x)'=\cos x$. Une périodicité se prouve par une égalité pour tout $x$.
\end{retenir}
\end{finchapitre}
\end{document}
